\subsection{n 次根}

设 \(n\) 是任一整数 \(> 0\) 。由关系 \(0 < x < y\) ，对 \(n\) 用归纳法推出 \(0 < x^n < y^n\) 。换句话说，函数 \(x \to x^n\) 对 \(x \geqslant 0\) 严格递增；它显然在每一点连续，因此（§ 2 第 6 目，定理 5）是 \(\mathbf{R}_{+}\) 到一个区间 \(\mathbf{I}\) 上的同胚。另一方面，由于 \(x \geqslant 1\) 蕴含 \(x^{n-1} \geqslant 1\) ，从而 \(x^n \geqslant x\) ，可见 \(\mathbf{I}\) 不是有界的，因而 \(\mathbf{I} = \mathbf{R}_{+}\) 。对于 \(x \geqslant 0\) ，映射 \(x \to x^n\) 的逆的值记作 \(x^{1/n}\) 或 \(\sqrt[n]{x}\) ，称为 \(x\) 的 \(1\) 次幂或 \(x\) 的 n 次根（当 \(n = 2\) , 3 时，我们说平方根、立方根；当 \(n = 2\) 时，我们用 \(\sqrt{x}\) 代替 \(\sqrt[2]{x}\) 书写）。正数 \(x^{1/n}\) 于是定义为方程的唯一正解

\[
y ^ {n} = x \quad (x \geqslant 0).\tag{2}
\]

特别地，我们看到存在实数 x 使得 \(x^{2}=2\) ，而没有有理数具有这一性质；这样我们重新得到有理直线 Q 不是完备空间这一事实。

映射 \(x \to x^{1/n}\) （ \(\mathbf{R}_+\) 到其自身的）严格递增且连续。由 (2) 我们有 \(0^{1/n}=0, 1^{1/n}=1\) ，还有

\[
(x y) ^ {1 / n} = x ^ {1 / n} y ^ {1 / n};\tag{3}
\]

因而 \(x \to x^{1/n}\) 是拓扑群 \(\mathbf{R}_{+}^{*}\) 的一个自同构。

在第五章 § 4 第 1 目中，我们将通过求出乘法群 \(\mathbf{R}_{+}^{*}\) 的所有自同构来推广这一结果。

